\section{Evaluation}
\label{sec:eval}

We ask five questions, in the order a reader should want them answered. Does a
second memory tier raise \emph{goodput} --- requests served within a latency
target --- rather than merely raw throughput? Does the placement policy
contribute, or would uniform offloading do as well? What does a host-tier
request cost, and how does that cost scale with the interconnect? Does the
capacity model of \S\ref{sec:model} predict what we measure? And where does the
technique stop paying?

\subsection{Setup}
\label{sec:eval:setup}

\paragraph{Hardware.} A100-80 over PCIe~4 for the cost decomposition and the
capacity mechanism; GH200-96 over NVLink-C2C for everything that depends on a
coherent interconnect. Link bandwidth is \emph{measured} on both rather than
quoted: the widely cited 900\,GB/s for C2C is bidirectional, a KV stream is
unidirectional, and an independent measurement of the same hardware reports
419\,GB/s host-to-device~\cite{pie}.

\paragraph{Exclusive access, for every absolute.} Each number we report as an
absolute is taken on an idle, exclusively allocated device. This is not
boilerplate: the same measurement on a contended node and on an idle one
differed by $9.4\times$ (\S\ref{sec:impl:stash}). Ratios between arms run back
to back are robust to contention; absolutes are not.

\paragraph{Models.} Qwen2.5-32B fp16 as the primary model. On a 96\,GB device
its weights leave $\Ckv \approx 22$--28\,GB, which is the regime
\S\ref{sec:model} identifies as the one where a second tier pays. Llama-3.1-8B
as a negative control, where the model predicts no benefit. We deliberately do
not evaluate a 70B-class model: its capacity-limited batch is about seven
requests, so a single request is 14\% and no result could separate an effect
from a rounding artifact.

\paragraph{Workloads.} ShareGPT and Alpaca under Poisson arrivals, via
\texttt{vllm bench serve}. Two properties are required rather than preferred.
Real traces are heavy-tailed in both prompt and generation length, which is the
variation a \emph{per-request} placement policy exists to exploit --- on a
workload of identical requests our placement rule assigns every request the
same score and degenerates to arrival order by construction, so measuring it
equal to arrival order there would say nothing about the policy. And open-loop
arrivals are needed because a closed loop of $N$ simultaneous requests is a
different experiment from offering $\lambda$ requests per second; it also
invites metrics that reward slowness (below). These are the workloads Pie
reports on~\cite{pie}, which makes our numbers comparable to theirs.

\paragraph{Metric.} \textbf{SLO-attaining goodput}: requests per second whose
per-output-token latency met a target. Raw throughput is secondary, p99 latency
and achieved interconnect bandwidth are reported alongside so that ``the
transfer hides'' is measured rather than inferred.

We do \emph{not} report mean concurrency. It averages over sampled steps, so an
arm that serves requests more slowly holds them longer and scores
\emph{higher}. It is not a goodput metric, and an earlier result of ours that
relied on it was retracted for exactly that reason.

\paragraph{Baselines.} Stock vLLM with all tiering disabled, which is also
Pie's baseline and makes the two sets of numbers commensurable. We verify
separately that our fork with the feature off matches the unmodified engine to
within 3\%; a baseline quietly slowed by our own additions would inflate every
ratio here. Uniform offload at a matched byte budget represents the placement
that Pie and DirectKV instantiate at their own granularities. Where a
comparison against a published figure is drawn, we state which engine version
and workload each side used rather than presenting ratios as directly
commensurable.

\subsection{Does a second tier raise goodput?}
\label{sec:eval:capacity}

The decisive experiment holds the model, the device, the workload and the
substrate fixed and varies only whether a second memory tier is used. We sweep
the arrival rate upward until goodput stops rising --- that is, until the
device is saturated --- and report goodput at a latency target anchored to the
baseline's own measured step time rather than to a round number, since a target
below an arm's step time is unsatisfiable by every arm and measures nothing.

Two properties of this design are worth stating in advance, because they are
what make a result interpretable. A target at the baseline's own floor is the
one a host-tier request is \emph{least} able to meet, since it adds per-step
latency; a very loose target admits everything and discriminates nothing. The
informative regime is between them. And below the applicability gate of
Eq.~\ref{eq:gate} the model predicts \emph{no} benefit at any rate, so the
tightest points of the sweep are a prediction we can be wrong about in public.

\emph{The capacity mechanism, measured independently of the economics.} On
A100 we can separate whether host-tier requests are \emph{additive} from
whether they pay. Against a fixed KV pool, host-tier requests raise peak
concurrency from 8 to 48 --- and 48 is a lower bound, being the entire offered
load.
% src: paper/data/preload_capacity.json @ f3df05bdb
KV is delivered synthetically in this experiment, so it establishes capacity
only and its generated tokens are not meaningful; goodput with real KV is
\S\ref{sec:eval:cost}.

That result also isolates why the system model matters. An earlier design
capped a host-tier request's block \emph{growth} while leaving its prefill
resident. It freed about 11\% of the request's footprint while still paying
full interconnect every step, and the served total crept toward the all-HBM
capacity from below without ever passing it, at any setting of the resident
cap. A request whose prefill stays resident is strictly worse than one left in
device memory, at any link speed --- which is what makes the disaggregated
decode-node assumption (\S\ref{sec:background:pd}) load-bearing rather than
convenient.

\subsection{Is the placement policy load-bearing?}
\label{sec:eval:ablation}

If uniform offloading at the same interconnect budget performs as well, the
policy is decoration and the contribution reduces to the analysis. The arms add
one mechanism each: uniform offload at a matched byte budget; our admission
controller choosing how many while arrival order chooses which; scored
placement; and migration.

\textbf{Matched bytes, not matched request count.} Capping how many requests
are host-tier does not control this comparison, because the policies select
requests of different context lengths and therefore move different volume ---
part of any difference would be a difference in offload ratio wearing the
appearance of a policy win. We count bytes in both directions per arm and treat
arms differing by more than $1.25\times$ in traffic as not comparable. A
parallel control matches the resident-cohort cap, because an arm with a
different cap is a different \emph{allocation} rather than a different policy;
on A100 that distinction accounted for most of an apparent 14.5\% difference.

The adaptive probe's verdict is reported as a result in its own right. On PCIe
it measured demotion as harmful and reverted it, which is the correct action on
a link that cannot hide the transfer. On a coherent link it should measure
demotion as beneficial and keep it. If it does not, the host path carries a
fixed per-step cost that survives a fast interconnect, and that is a finding we
would rather publish than omit.

\subsection{What a host-tier request costs}
\label{sec:eval:cost}

Sweeping the number of host-tier requests on an otherwise identical engine and
fitting
$\mathrm{ITL}(r) = \mathrm{ITL}_0 + \text{fixed} + r\cdot c$
separates a fixed per-step cost of using the host path at all from a per-request
cost that fails to overlap compute. The two extrapolate to completely different
answers on a faster link, which is why a single operating point cannot settle
the question.

\begin{table}[t]
\centering
\small
\begin{tabular}{lrr}
\toprule
host-tier requests & ITL (ms) & tok/s \\
\midrule
0 & 49.99 & 57.5 \\
1 & 65.66 & 54.7 \\
2 & 65.93 & 55.8 \\
3 & 74.98 & 57.0 \\
\bottomrule
\end{tabular}
\caption{Qwen2.5-32B, 2048 context, 12 offered, 8 resident, A100/PCIe. Fit:
$9.6$\,ms fixed $+$ $4.7$\,ms per request, so \textbf{78\% of the transfer
overlaps compute} --- the interconnect is not the binding cost at this scale.
Throughput recovers toward the offload-free arm as requests are added but does
not cross it within this sweep: three host-tier requests remain 0.9\% below it.}
\label{tab:ridercost}
\end{table}
% src: paper/data/hide_efficiency.json @ f3df05bdb

Two readings of this data were ours and both were wrong, in ways worth
recording. A single two-request point cannot distinguish a steep slope from a
large intercept; read that way it suggested that 89\% of the transfer failed to
overlap, and only the sweep corrected it. And a concurrency-over-latency model
predicted break-even at seven requests where the measured throughput column
crosses between three and four, because it ignores queueing --- the
offload-free arm leaves requests waiting, and host-tier requests shorten that
queue. We report the measured crossing and keep the model only as a bound.

The fixed-versus-slope split is also unstable at three points: two runs of the
same sweep fitted $12.3 + 4.5$ and $9.6 + 4.7$. Because the fixed term is
precisely the part that a faster interconnect does \emph{not} shrink, that
split determines break-even on a coherent link, so the interconnect comparison
uses ten host-tier counts rather than three.

\paragraph{Swap semantics versus streaming.} A design whose offloaded unit is
\emph{owned} by one side must write that unit back once per generated token; one
whose device-side copy is a \emph{cache} need not (\S\ref{sec:related}). We
measure the difference directly by writing each layer's slot back to the host
mirror with the policy otherwise unchanged. The emulation is exact in extent:
both arms read identical volume, and the write-back arm moves \textbf{2.00×}
the total link traffic, with identical token output.
% src: paper_runs/20260911_ablation_smoke.json

\subsection{Does the capacity model predict it?}
\label{sec:eval:model}

Two parameters of Eq.~\ref{eq:step} are re-measured on each platform: the
weight-read efficiency $\effw$ and the paged-KV read efficiency $\effkv$ against
working-set size. The second is the larger risk to our own projections and we
run it early for that reason --- $\effkv$ sits in the capacity denominator, so a
kernel generation that reads paged KV more efficiently shortens the step,
shrinks the transfer budget, and \emph{lowers} the predicted gain.

\textbf{One term of the law is asserted rather than measured, and we say
which.} The multiplier has two factors: a bandwidth ratio $\Xlink/\BW$ and an
amplifier $\HBM/\Ckv$. We validate it across \emph{bandwidth}, by running the
same system on two interconnect generations. The $\HBM/\Ckv$ dependence rests
on the parameter fit and on the model's agreement with a published system's own
results (\S\ref{sec:model:directkv}); isolating it cleanly would require one
model at two weight precisions, an experiment we chose not to run. A reader who
wants that axis measured is asking for a specific experiment, and this is the
honest account of the trade.

\subsection{When does it not pay?}
\label{sec:eval:negative}

\emph{Below the applicability gate.} Eq.~\ref{eq:gate} predicts no benefit when
the latency target is tighter than the full-pool step time; the tightest points
of the goodput sweep test that directly.

\emph{When the weights are small.} Llama-3.1-8B leaves so much device memory
after weights that there is nothing to buy. Our model predicts $\times$1.26 and
DirectKV independently reports that 8B does not benefit from offload. A model
that predicts its own failure cases is worth more than one tuned to its wins.

\emph{When the interconnect is slow.} On A100/PCIe, with real KV delivered
through the prefill handover, throughput is 2.3\% \emph{below} not offloading
at all.
% src: paper/data/pd_decode_goodput.json @ f3df05bdb
That is what the model predicts for a 25\,GB/s link, and DirectKV reaches the
same boundary from the kernel side, concluding that on PCIe theirs ``should be
viewed primarily as a capacity-extension mechanism\ldots bounded by CPU--GPU
interconnect bandwidth''. We report it as a validated negative that fixes one
end of the bandwidth axis rather than as a disappointment.

\subsection{Correctness}
\label{sec:eval:correct}

Placement moves bytes rather than approximating them, so these gates are
hygiene rather than a contribution. Numerically we verify the host pool's round
trip, the two-tier attention against a single-pass reference over the same KV
(576 mixed decode layer-steps, $\max|\mathrm{diff}| = 0$), and the prefill
handover across block boundaries and after the producer's memory has been
recycled. End to end we compare generated tokens against stock vLLM.

One methodological note, because it traps anyone reproducing this.
\texttt{bf16} attention reductions are not associative, so their result depends
on batch composition: stock vLLM produces different text for the same prompt at
concurrency 1 and at 6, reproducibly at each. Any system that changes batching
--- as a memory tier does --- will diverge on such prompts without being wrong.
We therefore compare serially, batch of one on both sides, and under load only
on prompts pre-screened as batching-stable. For calibration: on a serial greedy
comparison our host tier matched stock vLLM on 6 of 6 prompts in one round and
4 of 6 in another, while enabling \emph{vLLM's own prefix cache} --- changing
nothing else --- matched on 2 of 6.
